\section{Every zone price episode}
\label{app:episodes}

\begin{table}[H]
\centering
\caption{\textbf{All \zEpisodes{} episodes the detector finds, with their measured
response.} Weeks: Dominick's week numbers of the episode. Price: the store's
median log price against its zone peers over the episode, \%. Controls: same-zone
stores within 1\% of their zone. 8 weeks: difference-in-differences estimate on
log category units over the first eight weeks, \%, with its interval clustered
by store (too narrow, \S\ref{sec:zones}); Full: over the whole episode.
\textsuperscript{*} Listed, not scored: fewer than three controls, or
fewer than the 26 comparison weeks before the start (the four episodes starting
before week 27). An interval of zero width comes from an episode with one
control store, where the clustered standard error is zero, or, for soft drinks
store 105, from an episode with no comparison weeks.}
\label{tab:episodes}
\scriptsize\renewcommand{\arraystretch}{0.86}
\begin{tabular}{@{}llrrrrrrl@{}}
\toprule
Category & Zone & Store & Weeks & Price & Controls & 8 weeks & 95\% interval & Full \\
\midrule
Analgesics & 2 & 50 & 190--215 & 8.8 & 16 & 4.1 & 0.6 to 7.6 & −5.7 \\
Analgesics & 2 & 49 & 197--209 & 8.5 & 16 & −2.7 & −5.6 to 0.3 & −3.2 \\
Analgesics & 2 & 48 & 203--218 & 8.1 & 11 & −5.0 & −9.8 to −0.2 & −7.1 \\
Analgesics & 7 & 109 & 203--218 & 6.0 & 2 & 8.2 & −4.1 to 20.4 & 9.1 \textsuperscript{*} \\
Analgesics & 1 & 32 & 204--218 & −6.5 & 7 & 13.6 & 6.7 to 20.6 & 14.9 \\
Analgesics & 2 & 88 & 204--218 & −6.7 & 10 & −3.5 & −9.1 to 2.2 & −2.0 \\
Analgesics & 2 & 89 & 204--218 & 6.3 & 10 & −12.7 & −18.3 to −7.1 & −9.4 \\
Bottled juice & 12 & 105 & 9--24 & −7.9 & 7 & −7.2 & −18.9 to 4.4 & −3.2 \textsuperscript{*} \\
Bottled juice & 1 & 123 & 14--30 & −5.5 & 11 & 7.8 & 3.6 to 12.1 & 11.5 \textsuperscript{*} \\
Bottled juice & 6 & 40 & 17--28 & −6.6 & 6 & 8.2 & −9.0 to 25.3 & 5.8 \textsuperscript{*} \\
Bottled juice & 6 & 40 & 63--75 & −7.7 & 6 & −3.2 & −13.5 to 7.1 & −0.1 \\
Bottled juice & 2 & 84 & 266--283 & −9.4 & 24 & 11.0 & 8.9 to 13.2 & 12.7 \\
Bottled juice & 2 & 84 & 296--315 & −9.2 & 22 & −6.0 & −12.3 to 0.3 & −0.9 \\
Cheese & 6 & 77 & 249--261 & −7.1 & 7 & 4.8 & −12.3 to 21.9 & 4.1 \\
Frozen entrees & 5 & 8 & 201--218 & 10.2 & 1 & −26.8 & −26.8 to −26.8 & −31.0 \textsuperscript{*} \\
Frozen entrees & 2 & 47 & 202--218 & 6.7 & 9 & −4.7 & −10.1 to 0.7 & −3.2 \\
Frozen entrees & 2 & 28 & 204--218 & −6.9 & 9 & 8.0 & 2.4 to 13.5 & 7.9 \\
Frozen entrees & 1 & 71 & 204--218 & 9.3 & 4 & 21.6 & 12.7 to 30.6 & 17.1 \\
Frozen entrees & 2 & 91 & 204--218 & −8.4 & 9 & 6.6 & 1.0 to 12.1 & 9.7 \\
Frozen entrees & 6 & 40 & 205--218 & −7.6 & 2 & −3.2 & −6.8 to 0.4 & −2.1 \textsuperscript{*} \\
Frozen entrees & 12 & 114 & 205--218 & −8.7 & 3 & 9.2 & −5.7 to 24.0 & 7.5 \\
Frozen entrees & 2 & 119 & 205--218 & −8.0 & 9 & −0.1 & −5.6 to 5.4 & −1.1 \\
Frozen entrees & 1 & 2 & 220--237 & 9.0 & 4 & −6.2 & −11.0 to −1.5 & −12.8 \\
Frozen entrees & 5 & 8 & 220--237 & 9.8 & 2 & −13.0 & −95.9 to 69.9 & −7.7 \textsuperscript{*} \\
Frozen entrees & 2 & 9 & 220--237 & −9.6 & 9 & 13.3 & 9.4 to 17.2 & 7.8 \\
Frozen entrees & 7 & 33 & 220--237 & 5.9 & 2 & 2.3 & 1.0 to 3.6 & −1.1 \textsuperscript{*} \\
Frozen entrees & 6 & 40 & 220--237 & −10.9 & 1 & 21.0 & 21.0 to 21.0 & 15.6 \textsuperscript{*} \\
Frozen entrees & 2 & 49 & 220--234 & 8.9 & 9 & 4.7 & 0.8 to 8.6 & 5.0 \\
Frozen entrees & 1 & 68 & 220--237 & 7.6 & 4 & −3.1 & −7.8 to 1.7 & −12.2 \\
Frozen entrees & 1 & 71 & 220--237 & 10.0 & 4 & 7.4 & 2.7 to 12.2 & −5.4 \\
Frozen entrees & 7 & 75 & 220--234 & −10.2 & 2 & −0.3 & −1.6 to 1.0 & −4.7 \textsuperscript{*} \\
Frozen entrees & 2 & 76 & 220--234 & 8.6 & 9 & 4.7 & 0.8 to 8.6 & 4.8 \\
Frozen entrees & 1 & 95 & 220--234 & 8.1 & 4 & −1.7 & −6.4 to 3.1 & −6.4 \\
Frozen entrees & 2 & 110 & 220--234 & 8.4 & 9 & −3.7 & −7.6 to 0.2 & −1.6 \\
Frozen entrees & 2 & 119 & 220--237 & −6.6 & 9 & 2.0 & −1.9 to 5.9 & −2.8 \\
Frozen entrees & 1 & 130 & 220--233 & −7.3 & 4 & −3.3 & −8.0 to 1.5 & 2.1 \\
Frozen entrees & 1 & 71 & 270--290 & −7.0 & 14 & 9.4 & 6.0 to 12.8 & 5.0 \\
Soft drinks & 12 & 105 & 1--20 & −7.4 & 7 & 18.0 & 18.0 to 18.0 & 8.8 \textsuperscript{*} \\
Snack crackers & 2 & 50 & 196--216 & 9.6 & 10 & 9.4 & 3.2 to 15.6 & −13.1 \\
Snack crackers & 2 & 48 & 197--218 & 8.5 & 9 & −12.8 & −18.5 to −7.1 & −13.2 \\
Snack crackers & 2 & 76 & 197--218 & 6.9 & 9 & −11.9 & −18.1 to −5.8 & −6.8 \\
Snack crackers & 2 & 49 & 198--218 & 9.6 & 9 & −14.3 & −19.2 to −9.3 & −6.6 \\
Snack crackers & 2 & 89 & 198--218 & 8.5 & 9 & −7.7 & −12.6 to −2.8 & −1.1 \\
Snack crackers & 2 & 47 & 200--218 & 6.5 & 9 & 7.2 & 2.7 to 11.6 & 9.6 \\
Snack crackers & 2 & 28 & 203--218 & −6.5 & 9 & 24.0 & 20.2 to 27.8 & 22.4 \\
Snack crackers & 5 & 94 & 203--218 & −7.1 & 2 & −1.2 & −8.0 to 5.7 & −6.2 \textsuperscript{*} \\
Snack crackers & 1 & 2 & 204--218 & 8.7 & 4 & 5.5 & −7.4 to 18.3 & 2.9 \\
Snack crackers & 7 & 75 & 204--218 & −7.4 & 2 & −9.9 & −10.9 to −8.8 & −12.6 \textsuperscript{*} \\
Snack crackers & 2 & 88 & 204--218 & −8.6 & 9 & −2.9 & −7.2 to 1.3 & −0.6 \\
Snack crackers & 12 & 114 & 204--218 & −6.7 & 2 & 9.2 & −2.7 to 21.0 & 6.0 \textsuperscript{*} \\
Snack crackers & 2 & 119 & 204--218 & −6.5 & 9 & 3.5 & −0.8 to 7.7 & 4.1 \\
Snack crackers & 2 & 113 & 206--218 & −6.3 & 9 & 0.7 & −3.9 to 5.3 & 0.1 \\
Snack crackers & 6 & 40 & 207--218 & −8.7 & 2 & 9.5 & −18.8 to 37.8 & 14.6 \textsuperscript{*} \\
Snack crackers & 12 & 129 & 207--218 & −7.4 & 2 & 2.1 & −4.9 to 9.0 & −0.9 \textsuperscript{*} \\
\bottomrule
\end{tabular}

\end{table}

\section{The plan and its amendments}
\label{app:plan}

The plan (\texttt{docs/PAPER2\_PLAN.md}) was committed before any audit result
on 13 October 2026. It fixed the eight checks, the C7 detector, estimator and
per-episode score, and the rule that no audit result changes a model setting.
Eight dated entries followed. Entries 1 to 5 and 8 were made before the results
they affect were read; entries 6 and 7 were made after the earlier audit's
results, and everything they add is reported as added after the results.
\begin{enumerate}\setlength{\itemsep}{1pt}
\item 13 October, before any model lift: the episode detector measures each
store against its zone peers rather than the zone's most common price.
\item 13 October, after the episode list, before any model lift: at least three
controls to score an episode; models scored on an episode's first eight weeks,
against an 8-week difference in differences.
\item 13 October, before any audit result: the known-answer study, the C7
robustness analyses (event study, placebo, detector settings, pooled slope by
category), the pass thresholds for C1 to C6 and the main test.
\item 13 October, after the C7 robustness results, before any model lift for
C7: the price path and the pooled-slope intervals, labelled exploratory, and
the comparison of models with the measured response made on the pooled slope.
\item 13 October, correction, before any model lift for C7: an episode is
scored only if its 26 comparison weeks exist, leaving \zScored{} scored.
\item 17 October, after the audit results: the diagnosis of C1 and C4 on
fourteen saved runs; the main test also over versions; pass counts without C5,
and without C1 and C5, as exploratory.
\item 18 October, after the audit results: C1 under the model's own plan for
nothing planned; the comparison of the two plans for every configuration and
window; the benchmark described as observational; the main test rerun holding
the data and window fixed.
\item 13 November, before any audit of the tuned models: the audit moves to
each data set's final tuned world model; C1 scored under both plans; C3
between a model's copies; C6 against the tuned LightGBM; C7 and C8 with the
final settings; and the test of 18 October applied to the versions the
settings searches saved.
\end{enumerate}

\section{Every known-answer model}
\label{app:known}

\begin{table}[H]
\centering
\caption{\textbf{Every model in the known-answer study, with its explanation
scored against the truth.} Lift against truth: rank agreement and ratio of
summed lifts on item-weeks whose true lift is not zero, overall and for the
price, promotion and display parts. C4: item-weeks lower after a deeper cut. C5: summed forecast over
summed sales on planned weeks. Passed: checks passed of those that apply (C1 to
C6). RMSSE: item by store, weeks 253 to 260.}
\label{tab:known}
\footnotesize
\begin{tabular}{@{}lrrrrrrrrr@{}}
\toprule
 & \multicolumn{5}{c}{Lift against truth} & & & & \\
\cmidrule(lr){2-6}
Model & rank & sum & price & promotion & display & C4 & C5 & Passed & RMSSE \\
\midrule
\multicolumn{10}{@{}l}{\emph{Plans at random}} \\
World model, seed 1 & 0.99 & 1.14 & 1.13 & 1.08 & 1.00 & 0.0\% & 1.08 & 6/6 & 0.401 \\
World model, seed 2 & 0.99 & 1.10 & 1.09 & 1.06 & 1.02 & 0.0\% & 1.06 & 6/6 & 0.401 \\
World model, seed 3 & 0.99 & 1.11 & 1.11 & 1.04 & 0.91 & 0.0\% & 1.05 & 5/6 & 0.398 \\
World model, blind to price & 0.59 & 0.73 & 0.01 & 2.00 & 0.92 & 31.5\% & 0.90 & 3/5 & 0.510 \\
World model, blind to levers & 0.89 & 1.10 & 1.35 & 0.14 & 0.05 & 0.0\% & 1.09 & 4/5 & 0.450 \\
LightGBM & 0.85 & 1.09 & 1.03 & 0.98 & 1.39 & 7.6\% & 0.99 & 1/2 & 0.389 \\
LightGBM, blind to price & 0.43 & 0.73 & −0.00 & 1.95 & 1.41 & 49.5\% & 0.88 & 1/2 & 0.515 \\
\multicolumn{10}{@{}l}{\emph{Plans in high-demand weeks}} \\
World model, seed 1 & 0.96 & 1.33 & 1.22 & 1.29 & 1.83 & 0.0\% & 1.06 & 5/6 & 0.345 \\
World model, seed 2 & 0.95 & 1.35 & 1.24 & 1.31 & 1.81 & 0.0\% & 1.07 & 6/6 & 0.349 \\
World model, seed 3 & 0.95 & 1.40 & 1.27 & 1.38 & 1.93 & 0.0\% & 1.09 & 5/6 & 0.344 \\
World model, blind to price & 0.51 & 1.09 & −0.00 & 2.37 & 2.05 & 27.2\% & 0.93 & 2/5 & 0.489 \\
World model, blind to levers & 0.82 & 1.23 & 1.63 & 0.06 & −0.03 & 0.0\% & 1.07 & 4/5 & 0.480 \\
LightGBM & 0.82 & 1.26 & 1.13 & 1.07 & 1.63 & 9.9\% & 1.01 & 1/2 & 0.322 \\
LightGBM, blind to price & 0.54 & 0.94 & −0.02 & 2.02 & 1.76 & 43.4\% & 0.90 & 1/2 & 0.497 \\
\bottomrule
\end{tabular}

\end{table}
